\documentclass[11pt,a4paper]{article}
\usepackage[margin=25mm]{geometry}
\usepackage{amsmath,amssymb,amsthm,mathtools,microtype}
\usepackage[colorlinks=true,linkcolor=blue,urlcolor=blue,citecolor=blue]{hyperref}
\newtheorem{theorem}{Theorem}
\newtheorem{proposition}{Proposition}
\newcommand{\br}[2]{\{#1,#2\}}
\newcommand{\TT}{\mathbb T^2}
\title{\textbf{Algebraic Whistler Spectral Barriers and Nonlinear Chiral-Mode Leakage}\\in Two-and-a-Half-Dimensional Electron Magnetohydrodynamics}
\author{Tony Newton\\Newton Astro Labs\thanks{Newton Astro Labs is the trading name under which Tony Newton conducts independent computational research in the United Kingdom; it is not a limited company.}\\London, UK\\\texttt{tony.newton79@gmail.com}}
\date{}
\begin{document}
\maketitle
\begin{abstract}
We obtain a restricted spectral obstruction, an exact nonlinear chiral-mode coupling formula and a family of positive controls for ideal two-and-a-half-dimensional electron magnetohydrodynamics (EMHD). Its linearized Fourier generator has irreducible quadratic characteristic polynomial over \(\mathbb C(k_x,k_y)\); rational-symbol scalar diagonalization is therefore impossible, and a square-root nonlocal Fourier multiplier is necessary and sufficient. An exact quadratic coefficient measures instantaneous nonlinear leakage between chiral whistler families: a pair of noncollinear planar modes leaks precisely when their wavenumber magnitudes differ. Modes of equal magnitude, by contrast, generate finite-amplitude exact Beltrami-shell solutions. The spectral and interaction statements are restricted to their indicated algebraic classes; they do not address nonlinear Lax integrability. The elementary point-potential flatness calculation is a special case of established zero-curvature theory, not a separate contribution. We neither construct nor exclude non-removable-parameter Lax representations for nonlinear EMHD. Priority of the narrower spectral and convolution formulas remains unresolved.
\end{abstract}

\section{Model and assumptions}
Define \(\{f,g\}=f_xg_y-f_yg_x\), \(\Delta=\partial_x^2+\partial_y^2\).
With magnetic field \(\mathbf B=(a_y,-a_x,b)\), ideal inertia-free EMHD reduces to
\begin{equation}\label{eq:emhd}
 a_t=\{a,b\},\qquad b_t=\{\Delta a,a\}.
\end{equation}
Let \(a=B_0y+\psi\), \(b=\beta\), where \(B_0\ne0\), and let the perturbations be smooth and \(2\pi\)-periodic in \(x,y\). The physical magnetic field is periodic even though its background potential has an affine term. The exact perturbation equations read
\begin{equation}\label{eq:pert}
 \psi_t=-B_0\beta_x+\{\psi,\beta\},\qquad
 \beta_t=B_0(\Delta\psi)_x+\{\Delta\psi,\psi\}.
\end{equation}
Let \(\Lambda=\sqrt{-\Delta}\) be the periodic multiplier \(\widehat{\Lambda f}(k)=|k|\widehat f(k)\).
The 2.5D reduction and nonlinear whistlers are established subjects
\cite{peng,lyutikov,shaikh}. Morozov's representations for reduced and two-dimensional ideal MHD
\cite{morozov} concern equations distinct from \eqref{eq:emhd}, but his
covering and spectral methods demonstrate why a point-potential restriction
cannot settle the existence of more general nonlinear EMHD Lax representations.
Marvan's foundational zero-curvature theory \cite{marvan1993,marvan2002}
provides the relevant setting for gauge and parameter equivalence.

\section{Minimal spectral algebraic extension}
The linearized Fourier generator at \(k=(k_x,k_y)\) is
\begin{equation}\label{eq:matrix}
 A(k)=-iB_0k_x
 \begin{pmatrix}0&1\\k_x^2+k_y^2&0\end{pmatrix},\qquad
 \det(zI-A)=z^2+B_0^2k_x^2(k_x^2+k_y^2).
\end{equation}
\begin{theorem}[Sharp rational diagonalization barrier]\label{thm:barrier}
Let \(K=\mathbb C(k_x,k_y)\), \(B_0\ne0\).
The characteristic polynomial of \eqref{eq:matrix} is irreducible in \(K[z]\). Therefore no invertible \(K\)-valued similarity transformation can diagonalize \(A\) into scalar rational Fourier symbols. Nor can a one-field first-order-in-time evolution with rational Fourier symbol produce either whistler eigenfrequency on generic planar modes. The smallest splitting field is \(K(s)\) with \(s^2=k_x^2+k_y^2\). In physical variables, \(w_\pm=\beta\pm\Lambda\psi\) satisfy exactly
\begin{equation}\label{eq:diag}
 (w_+)_t=-B_0\Lambda\partial_xw_+,\qquad
 (w_-)_t=+B_0\Lambda\partial_xw_-,
\end{equation}
under the linearized equation.
\end{theorem}
\begin{proof}
Over \(\mathbb C[k_x,k_y]\),
\(k_x^2+k_y^2=(k_x+ik_y)(k_x-ik_y)\).
The two distinct irreducible linear factors occur with odd exponents, hence this polynomial is not a square in \(K\). Multiplication by \((iB_0k_x)^2\) preserves its square class. Thus the quadratic characteristic polynomial has no root in \(K\). A rational-symbol diagonalizer would produce such eigenvalues in \(K\), contradiction. Both eigenvalues are in the quadratic extension \(K(s)\), which is minimal. For \(s\ne0\), the matrix
\[
 S=\begin{pmatrix}s&1\\-s&1\end{pmatrix}
\]
satisfies \(SAS^{-1}=\operatorname{diag}(-iB_0k_xs,+iB_0k_xs)\). Substituting \(s\mapsto\Lambda\) proves \eqref{eq:diag}.
\end{proof}
This excludes scalar \emph{rational} diagonalization, not local matrix evolution: the original \(2\times2\) system is local and already realizes both branches. Such a linear evolution matrix is not in itself a Lax pair.

\section{Exact nonlinear chiral leakage}
Using \(w_\pm=\beta\pm\Lambda\psi\), the nonlinear system \eqref{eq:pert} is exactly
\begin{equation}\label{eq:chiral}
 (w_\pm)_t=\mp B_0\Lambda\partial_xw_\pm+
 \{\Delta\psi,\psi\}\pm\Lambda\{\psi,\beta\}.
\end{equation}
On initial data with \(w_-=0\), hence \(\beta=\Lambda\psi\), the normal time derivative is
\begin{equation}\label{eq:leak}
 (w_-)_t=\mathcal D(\psi):=\{\Delta\psi,\psi\}-\Lambda\{\psi,\Lambda\psi\}.
\end{equation}
\begin{theorem}[Pairwise chiral-leakage multiplier]\label{thm:pair}
For any periodic trigonometric polynomial \(\psi=\sum_p\psi_pe^{ip\cdot x}\),
\begin{equation}\label{eq:fourier}
 \widehat{\mathcal D(\psi)}(k)=\frac12
 \sum_{p+q=k}(p_xq_y-p_yq_x)
 (|p|-|q|)(|p|+|q|-|p+q|)\psi_p\psi_q.
\end{equation}
For two distinct nonzero modes \(\psi=Ae^{ip\cdot x}+Ce^{iq\cdot x}\) with \(AC\ne0\), the coefficient at \(p+q\) is
\begin{equation}\label{eq:two}
 AC(p_xq_y-p_yq_x)(|p|-|q|)(|p|+|q|-|p+q|).
\end{equation}
It is nonzero exactly when the wavevectors are noncollinear and have different lengths.
\end{theorem}
\begin{proof}
The Fourier bracket formula is
\(\widehat{\{f,g\}}(k)=-\sum_{p+q=k}(p\times q)\widehat f(p)\widehat g(q)\).
Substituting \(\widehat{\Delta\psi}(p)=-|p|^2\psi_p\) and \(\widehat{\Lambda\psi}(q)=|q|\psi_q\) in \eqref{eq:leak} gives the unsymmetrized coefficient
\(\sum_{p+q=k}(p\times q)(|p|^2+|k||q|)\psi_p\psi_q\).
Average with the same sum after swapping \(p,q\). Its coefficient becomes
\(\tfrac12(p\times q)(|p|-|q|)(|p|+|q|-|k|)\), proving \eqref{eq:fourier}. For two modes the ordered pairs add to \eqref{eq:two}; self-pairs vanish. For noncollinear vectors, both \(p\times q\ne0\) and the strict triangle inequality \(|p+q|<|p|+|q|\) hold. The only remaining zero condition is \(|p|=|q|\).
\end{proof}
A real exact witness is
\[
 \psi=\cos x+\cos(2y),\qquad
 \beta=\cos x+2\cos(2y)=\Lambda\psi.
\]
Here \(\{\Delta\psi,\psi\}=6\sin x\sin(2y)\), \(\{\psi,\Lambda\psi\}=2\sin x\sin(2y)\), and \(\Lambda\sin x\sin(2y)=\sqrt5\sin x\sin(2y)\), so
\begin{equation}\label{eq:witness}
 (w_-)_t=(6-2\sqrt5)\sin x\sin(2y)\ne0
\end{equation}
at the initial data, independently of \(B_0\). This formal derivative does not invoke an ideal-EMHD Cauchy existence result. If a twice time-differentiable smooth solution exists through these data, then \(E_-=\frac14\int_{\mathbb T^2}w_-^2\) obeys
\(E_-'(0)=0\) and \(E_-''(0)=\frac{\pi^2}{2}(6-2\sqrt5)^2>0\).

\section{Exact finite-amplitude equal-radius shells}
\begin{proposition}[Beltrami whistler positive control]\label{prop:wave}
If \(\kappa>0\), \(\sigma=\pm1\), and
\(-\Delta\psi_0=\kappa^2\psi_0\), then
\begin{equation}\label{eq:wave}
 \psi(x,y,t)=\psi_0(x-\sigma B_0\kappa t,y),\qquad
 \beta(x,y,t)=\sigma\kappa\psi(x,y,t)
\end{equation}
is an exact finite-amplitude solution of both nonlinear equations \eqref{eq:pert}. The eigenfunction may superpose noncollinear wavevectors on the same Fourier circle.
\end{proposition}
\begin{proof}
At every time \(\Delta\psi=-\kappa^2\psi\), \(\beta=\sigma\kappa\psi\), so \(\{\psi,\beta\}=\{\Delta\psi,\psi\}=0\). The surviving equations are
\(\psi_t=-\sigma B_0\kappa\psi_x\) and
\(\beta_t=-B_0\kappa^2\psi_x\), exactly \eqref{eq:wave}. The magnetic perturbation \((\psi_y,-\psi_x,\beta)\) is Beltrami:
\(\nabla\times\delta B=\sigma\kappa\delta B\).
\end{proof}
This construction is included as a known-type harmonic/Beltrami whistler control, not claimed to be new \cite{lyutikov}.
\begin{proposition}[Energy consistency]
Any sufficiently smooth periodic solution conserves
\[
 E=\frac12\int_{\mathbb T^2}(\beta^2+|\nabla\psi|^2)
 =\frac14\int_{\mathbb T^2}(w_+^2+w_-^2).
\]
\end{proposition}
\begin{proof}
Let \(q=\Delta\psi\). The derivative is
\[
 \dot E=\int_{\mathbb T^2}\bigl[
 \beta(B_0q_x+\{q,\psi\})-q(-B_0\beta_x+\{\psi,\beta\})\bigr].
\]
The background terms are a total derivative; the Jacobian terms cancel by cyclic integration. The final formula follows from Parseval.
\end{proof}

\section{Relation to zero-curvature representations}
The irreducibility result in Section~2 concerns scalar rational diagonalization
of the \emph{linearized} Fourier generator. It is not an obstruction
to a spectral parameter in a nonlinear Lax representation. The relevant
notions of zero-curvature representation, gauge equivalence and parameter
removability are developed in the general framework of Marvan
\cite{marvan1993,marvan2002}; the differential-covering methods used by
Morozov \cite{morozov} lie outside our rational-symbol calculation.

For completeness, consider the elementary restriction to a matrix
connection
\[
 \Phi_x=U\Phi,\qquad \Phi_y=V\Phi,\qquad \Phi_t=W\Phi,
\]
where the potentials depend pointwise only on
\((\psi,\beta,q)\), \(q=\Delta\psi\), and possibly an auxiliary
parameter. If the connection is flat for every freely specified
spatial solution jet, comparison of the independent spatial derivatives
in the \(xy\)-curvature forces \(U\) and \(V\) to be independent of
\(\psi,\beta,q\), with \([U,V]=0\). The \(xt\)- and \(yt\)-curvatures
then force \(W\) to be field-independent and all three matrices to
commute. This is immediate jet coefficient matching in a restrictive
class, rather than an independent integrability theorem.

Derivative-dependent matrix connections, differential and
pseudodifferential Lax operators, and nonlocal coverings are not
subject to this conclusion. In particular, the present calculation
neither constructs nor rules out non-removable-parameter Lax
representations for the full nonlinear EMHD system \eqref{eq:emhd}.

\section{Verification, limits and next question}
The bundled exact SymPy verifier checks the characteristic polynomial, algebraic diagonalization, two-mode coupling coefficient, real leakage witness and both helicities of a noncollinear equal-radius wave. A separate standard-library-only verifier independently checks integer mode-pair geometry and the strict positivity of the real leakage coefficient. Their coverage is explicitly stated in the accompanying audit. These checks support, but do not replace, the proofs above.

The square-root diagonalization has familiar wave-physics antecedents; whistler interactions and exact Beltrami waves are known \cite{lyutikov}. Morozov's representations \cite{morozov} concern related but distinct MHD equations, while the jet-space and gauge-equivalence methods of Marvan \cite{marvan1993,marvan2002} place the restricted point-potential calculation in established prior art. We have not established priority of the particular spectral or convolution identities. No nonlinear EMHD Lax classification, proof of spectral-parameter non-removability, general nonintegrability theorem, or new plasma observation follows from the present results.

\begin{thebibliography}{9}
\bibitem{peng} Q. Peng, \emph{Local well-posedness for the two-and-a-half-dimensional EMHD system with split fractional dissipation}, arXiv:2605.20845 (2026), \url{https://arxiv.org/abs/2605.20845}.
\bibitem{lyutikov} M. Lyutikov, \emph{Electron magnetohydrodynamics: Dynamics and turbulence}, Phys. Rev. E \textbf{88}, 053103 (2013), \url{https://doi.org/10.1103/PhysRevE.88.053103}.
\bibitem{morozov} O. I. Morozov, \emph{Lax representations and variational Poisson structures for magnetohydrodynamics equations}, Analysis and Mathematical Physics \textbf{15}, 120 (2025), \url{https://doi.org/10.1007/s13324-025-01119-w}.
\bibitem{marvan1993} M. Marvan, \emph{On zero-curvature representations of partial differential equations}, in \emph{Differential Geometry and Its Applications}, Proceedings of the Opava Conference (1992), Mathematical Publications \textbf{1}, Silesian University, Opava (1993), pp.~103--122.
\bibitem{marvan2002} M. Marvan, \emph{On the horizontal gauge cohomology and nonremovability of the spectral parameter}, Acta Applicandae Mathematicae \textbf{72}, 51--65 (2002), \url{https://doi.org/10.1023/A:1015218422059}.
\bibitem{shaikh} D. Shaikh, \emph{Theory and simulations of whistler wave propagation}, J. Plasma Phys. \textbf{75}, 117--132 (2009), \url{https://doi.org/10.1017/S0022377808007198}.
\end{thebibliography}
\end{document}
